% =====================================================================
\subsection{Solutions to Form A}\label{sec:sol-a}
% =====================================================================

Each solution is written in full, the way it should appear on an exam paper.
Point allocations are shown in square brackets. Faster routes, when they
exist, are shown only afterwards in an \emph{Alternative method} box.

% ---------------------------------------------------------------------
\solhead{1}{5}
\textbf{Idea.} A vector function is defined exactly where \emph{all} of its
component functions are defined, so the domain is the intersection of the three
component domains.
\begin{itemize}
\item $\sqrt{4-t^2}$ requires $4-t^2\ge0$, i.e. $t^2\le4$, i.e. $-2\le t\le2$.
\item $\ln(t+1)$ requires $t+1>0$, i.e. $t>-1$ (the logarithm needs a \emph{positive} argument).
\item $\dfrac{1}{t-1}$ requires $t-1\neq0$, i.e. $t\neq1$.
\end{itemize}
Intersecting: from $[-2,2]$ and $t>-1$ we get $(-1,2]$; removing $t=1$ leaves
\[
  \boxed{(-1,1)\cup(1,2]}.
\]
\textit{[2 pts for the three conditions, 3 pts for the correct intersection.]}

% ---------------------------------------------------------------------
\solhead{2}{10}
\textbf{(a)} The cylinder $x^2+y^2=4$ is a circle of radius $2$ when seen from above,
so we parametrize it by $x=2\cos t$, $y=2\sin t$, $0\le t\le2\pi$. Every point of the
curve also lies on the plane $x+z=3$, so $z=3-x=3-2\cos t$. Therefore
\[
  \boxed{\vr(t)=\ip{2\cos t,\;2\sin t,\;3-2\cos t},\qquad0\le t\le2\pi.}
\]
Check: $x^2+y^2=4\cos^2t+4\sin^2t=4$\checkok\ and $x+z=2\cos t+3-2\cos t=3$\checkok.

\textbf{(b)} \step{1}{Parameter of the point $(0,2,3)$.} We need $2\cos t=0$ and
$2\sin t=2$, i.e. $\cos t=0$ and $\sin t=1$, so $t=\frac\pi2$. The third coordinate
agrees: $3-2\cos\frac\pi2=3$.\checkok

\step{2}{Tangent vector.} $\vr\,'(t)=\ip{-2\sin t,\;2\cos t,\;2\sin t}$ (the derivative of
$3-2\cos t$ is $2\sin t$). At $t=\frac\pi2$: $\vr\,'\!\left(\frac\pi2\right)=\ip{-2,\;0,\;2}$.

\step{3}{Tangent line.} The line through $(0,2,3)$ in the direction $\ip{-2,0,2}$ has
parametric equations
\[
  \boxed{x=0-2s=-2s,\qquad y=2+0\cdot s=2,\qquad z=3+2s,\qquad s\in\R .}
\]
\textit{[(a) 5 pts, (b) 5 pts.]}

% ---------------------------------------------------------------------
\solhead{3}{5}
\textbf{Idea.} Integrate each component, then use the initial condition to find the
constant vector.

\step{1}{Integrate.} $\int2t\,\dd t=t^2$, $\int3t^2\,\dd t=t^3$, $\int e^t\,\dd t=e^t$. So
\[
  \vr(t)=\ip{t^2,\;t^3,\;e^t}+\vb C,
\]
where $\vb C=\ip{C_1,C_2,C_3}$ is a constant \emph{vector} (one constant per component).

\step{2}{Initial condition.} At $t=0$: $\vr(0)=\ip{0,0,e^0}+\vb C=\ip{0,0,1}+\vb C$. This must equal
$\ip{1,0,-1}$, so $\vb C=\ip{1,0,-1}-\ip{0,0,1}=\ip{1,0,-2}$.

\step{3}{Answer.}
\[
  \boxed{\vr(t)=\ip{t^2+1,\;t^3,\;e^t-2}}.
\]
Check: $\vr(0)=\ip{1,0,-1}$\checkok\ and $\vr\,'(t)=\ip{2t,3t^2,e^t}$\checkok.

% ---------------------------------------------------------------------
\solhead{4}{7}
\textbf{Idea.} $L=\int_1^e\norm{\vr\,'(t)}\,\dd t$.

\step{1}{Velocity.} $\vr\,'(t)=\ip{2,\;2t,\;\frac1t}$ (since $\frac{\dd}{\dd t}\ln t=\frac1t$).

\step{2}{Speed.}
\[
  \norm{\vr\,'(t)}=\sqrt{2^2+(2t)^2+\left(\tfrac1t\right)^2}=\sqrt{4+4t^2+\frac{1}{t^2}} .
\]
To take this square root, notice that the expression under it is a perfect square:
\[
  \left(2t+\frac1t\right)^2=(2t)^2+2\cdot2t\cdot\frac1t+\left(\frac1t\right)^2=4t^2+4+\frac{1}{t^2}.
\]
Since $2t+\frac1t>0$ for $t$ in $[1,e]$, the square root is $\norm{\vr\,'(t)}=2t+\dfrac1t$.

\step{3}{Integrate.}
\[
  L=\int_1^e\left(2t+\frac1t\right)\dd t=\Big[t^2+\ln t\Big]_1^e=\left(e^2+\ln e\right)-\left(1+\ln1\right)
  =(e^2+1)-(1+0)=\boxed{e^2}.
\]
\textit{[Speed 2 pts, recognizing the perfect square 3 pts, integration 2 pts.]}

\trap{Trying to integrate $\sqrt{4+4t^2+1/t^2}$ directly. Whenever an arc-length
integrand looks impossible, look for a perfect square of the form $(A+B)^2=A^2+2AB+B^2$.}

% ---------------------------------------------------------------------
\solhead{5}{8}
\textbf{Idea.} $\kappa=\dfrac{\norm{\vr\,'\times\vr\,''}}{\norm{\vr\,'}^3}$ at the parameter of the point.

\step{1}{Parameter.} $(t,t^2,t^3)=(1,1,1)$ when $t=1$.
\step{2}{Derivatives.} $\vr\,'=\ip{1,2t,3t^2}$ and $\vr\,''=\ip{0,2,6t}$; at $t=1$,
$\vr\,'(1)=\ip{1,2,3}$ and $\vr\,''(1)=\ip{0,2,6}$.
\step{3}{Cross product.}
\[
  \vr\,'(1)\times\vr\,''(1)=\begin{vmatrix}\vi&\vj&\vk\\1&2&3\\0&2&6\end{vmatrix}
  =\vi(2\cdot6-3\cdot2)-\vj(1\cdot6-3\cdot0)+\vk(1\cdot2-2\cdot0)=\ip{6,\;-6,\;2}.
\]
\step{4}{Lengths.} $\norm{\ip{6,-6,2}}=\sqrt{36+36+4}=\sqrt{76}=2\sqrt{19}$ and
$\norm{\vr\,'(1)}=\sqrt{1+4+9}=\sqrt{14}$, so $\norm{\vr\,'(1)}^3=14\sqrt{14}$.
\step{5}{Curvature.}
\[
  \kappa(1)=\frac{2\sqrt{19}}{14\sqrt{14}}=\frac{\sqrt{19}}{7\sqrt{14}}
  =\frac{\sqrt{19}\cdot\sqrt{14}}{7\cdot14}=\boxed{\frac{\sqrt{266}}{98}}\approx0.166 .
\]

% ---------------------------------------------------------------------
\solhead{6}{8}
\step{1}{Partial derivatives.} $f_x=3x^2-2y^2$ and $f_y=-4xy$. They are polynomials, hence
continuous, so $f$ is differentiable and $\Du f=\grad f\cdot\vu$.
\step{2}{Gradient at $P(1,-1)$.} $f_x=3(1)-2(1)=1$ and $f_y=-4(1)(-1)=4$, so $\grad f(1,-1)=\ip{1,4}$.
\step{3}{Unit vector toward $Q$.} $\vecPQ=\ip{4-1,\;3-(-1)}=\ip{3,4}$, length $5$, so $\vu=\ip{\frac35,\frac45}$.
\step{4}{Directional derivative.}
\[
  \Du f(1,-1)=1\cdot\frac35+4\cdot\frac45=\frac35+\frac{16}{5}=\boxed{\frac{19}{5}}.
\]

% ---------------------------------------------------------------------
\solhead{7}{15}
\step{1}{Gradient.} $T_x=2x$, $T_y=2y$, $T_z=-4z$, so $\grad T(1,1,1)=\ip{2,\;2,\;-4}$.

\textbf{(a)} The direction from $P(1,1,1)$ toward the origin is $\ip{0-1,0-1,0-1}=\ip{-1,-1,-1}$,
of length $\sqrt3$, so $\vu=\frac{1}{\sqrt3}\ip{-1,-1,-1}$. Then
\[
  \Du T(P)=\frac{1}{\sqrt3}\big(2(-1)+2(-1)+(-4)(-1)\big)=\frac{-2-2+4}{\sqrt3}=\boxed{0}.
\]
This is not an accident: $T(1,1,1)=1+1-2=0$, and the whole segment from $P$ to the
origin lies on the level surface $T=0$ (for points $(s,s,s)$, $T=s^2+s^2-2s^2=0$). Along a level
surface the function does not change, so the rate is $0$.

\textbf{(b)} $T$ increases most rapidly in the direction of the gradient $\ip{2,2,-4}$, i.e. of
the unit vector $\frac{1}{\sqrt{24}}\ip{2,2,-4}=\frac{1}{\sqrt6}\ip{1,1,-2}$. The maximum rate is
$\norm{\ip{2,2,-4}}=\sqrt{4+4+16}=\sqrt{24}=\boxed{2\sqrt6}$.

\textbf{(c)} The level surface through $P$ is $T=T(P)=0$, i.e. the cone $x^2+y^2-2z^2=0$.
The gradient $\ip{2,2,-4}$ is normal to it at $P$, so the tangent plane is
\[
  2(x-1)+2(y-1)-4(z-1)=0\iff2x+2y-4z=0\iff\boxed{x+y-2z=0}.
\]
\textit{[(a) 5 pts, (b) 5 pts, (c) 5 pts.]}

% ---------------------------------------------------------------------
\solhead{8}{12}
\step{1}{Gradient.} $f_x=2x+y$, $f_y=x$, so $\grad f(1,1)=\ip{3,1}$, with length $\sqrt{10}\approx3.16$.
\step{2}{Equations for $\vu=\ip{a,b}$.} Directional derivative $1$: $3a+b=1$. Unit length: $a^2+b^2=1$.
\step{3}{Solve.} From the first, $b=1-3a$. Substituting,
\[
  a^2+(1-3a)^2=1\;\Longrightarrow\;a^2+1-6a+9a^2=1\;\Longrightarrow\;10a^2-6a=0\;\Longrightarrow\;2a(5a-3)=0 .
\]
\begin{itemize}
\item $a=0\Rightarrow b=1$: $\vu=\ip{0,1}$.
\item $a=\frac35\Rightarrow b=1-\frac95=-\frac45$: $\vu=\ip{\frac35,-\frac45}$.
\end{itemize}
Check: $3\cdot0+1=1$\checkok; $3\cdot\frac35-\frac45=\frac95-\frac45=1$\checkok.
\[
  \boxed{\vu=\ip{0,1}\quad\text{or}\quad\vu=\ip{\tfrac35,-\tfrac45}}
\]
\step{4}{Can the rate be $4$?} No. For every unit vector,
$\Du f=\norm{\grad f}\cos\theta\le\norm{\grad f}=\sqrt{10}\approx3.16$, and $4>\sqrt{10}$. (Algebraically,
$3a+b=4$ together with $a^2+b^2=1$ has no real solution.)

% ---------------------------------------------------------------------
\solhead{9}{6}
\step{1}{Compute the vectors.} $\vF(x,y)=\ip{1,x}$ depends only on $x$:
\begin{itemize}
\item on the line $x=-1$ (at $y=-1,0,1$): $\vF=\ip{1,-1}$;
\item on the line $x=0$: $\vF=\ip{1,0}$;
\item on the line $x=1$: $\vF=\ip{1,1}$.
\end{itemize}
\step{2}{Sketch.} Draw each vector starting at its point (here at half length):
\begin{center}
\begin{tikzpicture}[scale=1.1,>=Stealth]
  \draw[->,gray] (-2,0) -- (2.6,0) node[right]{$x$};
  \draw[->,gray] (0,-1.8) -- (0,1.9) node[above]{$y$};
  \foreach \x in {-1,0,1}{ \foreach \y in {-1,0,1}{
     \fill[NavyBlue] (\x,\y) circle (1.5pt);
     \draw[->,thick,BrickRed] (\x,\y) -- ({\x+0.5},{\y+0.5*\x});
  }}
  \node[below left] at (0,0) {\small $0$};
\end{tikzpicture}
\end{center}
\step{3}{Describe.} Every arrow points to the right (first component $1$). The tilt is
determined by $x$: arrows point down-right for $x<0$, straight right on the $y$-axis,
and up-right for $x>0$, tilting more steeply as $\abs x$ grows. All arrows on the same
vertical line are identical, because $\vF$ does not depend on $y$.

% ---------------------------------------------------------------------
\solhead{10}{8}
\step{1}{Component test.} $M=2x+y\cos(xy)$, $N=x\cos(xy)+3y^2$.
By the product and chain rules,
\[
  M_y=0+\cos(xy)+y\cdot(-\sin(xy))\cdot x=\cos(xy)-xy\sin(xy),\qquad
  N_x=\cos(xy)+x\cdot(-\sin(xy))\cdot y+0=\cos(xy)-xy\sin(xy).
\]
These are equal and continuous on $\R^2$, so $\vF$ \textbf{is conservative}.
\step{2}{Integrate $M$ in $x$.} $\int2x\,\dd x=x^2$ and $\int y\cos(xy)\,\dd x=\sin(xy)$
(check: $\pd{}{x}\sin(xy)=y\cos(xy)$). So $f=x^2+\sin(xy)+g(y)$.
\step{3}{Match $N$.} $f_y=x\cos(xy)+g'(y)$ must equal $x\cos(xy)+3y^2$, so $g'(y)=3y^2$ and $g(y)=y^3+K$.
\[
  \boxed{f(x,y)=x^2+\sin(xy)+y^3+K}.
\]
\step{4}{Check.} $f_x=2x+y\cos(xy)=M$\checkok, $f_y=x\cos(xy)+3y^2=N$\checkok.

% ---------------------------------------------------------------------
\solhead{11}{9}
\step{1}{Component test.} $M=2xz+y^2$, $N=2xy$, $P=x^2+3z^2$:
$M_y=2y=N_x$;\ $N_z=0=P_y$;\ $P_x=2x=M_z$. All three hold on $\R^3$, so $\vF$ is conservative.
\step{2}{Integrate $M$ in $x$.} $f=x^2z+xy^2+g(y,z)$.
\step{3}{Match $N$.} $f_y=2xy+g_y=2xy\Rightarrow g_y=0\Rightarrow g=h(z)$, a function of $z$ only.
\step{4}{Match $P$.} $f_z=x^2+h'(z)=x^2+3z^2\Rightarrow h'(z)=3z^2\Rightarrow h(z)=z^3+K$.
So $f=x^2z+xy^2+z^3+K$.
\step{5}{Use $f(0,0,0)=5$.} $f(0,0,0)=K=5$.
\[
  \boxed{f(x,y,z)=x^2z+xy^2+z^3+5}.
\]

% ---------------------------------------------------------------------
\solhead{12}{7}
\step{1}{Write the conditions.} $M=axy-z^3$, $N=(a-2)x^2$, $P=(1-a)xz^2$.
\begin{itemize}
\item $M_y=ax$ and $N_x=2(a-2)x$. Equal for all $x$ iff $a=2(a-2)=2a-4$, i.e. $a=4$.
\item $N_z=0$ and $P_y=0$: always equal.
\item $P_x=(1-a)z^2$ and $M_z=-3z^2$. Equal iff $1-a=-3$, i.e. $a=4$.
\end{itemize}
Both nontrivial conditions give the same value, so $\boxed{a=4}$, and then
$\vF=\ip{4xy-z^3,\;2x^2,\;-3xz^2}$.
\step{2}{Potential.} $f=\int(4xy-z^3)\,\dd x=2x^2y-xz^3+g(y,z)$.
$f_y=2x^2+g_y=2x^2\Rightarrow g=h(z)$. $f_z=-3xz^2+h'(z)=-3xz^2\Rightarrow h=K$.
\[
  \boxed{f(x,y,z)=2x^2y-xz^3+K}.
\]
